\section{In-Process Object}

An In-Process object represents the intermediate state of the workpiece during a machining process.

In-process objects are optional and may be generated by simulation, verification, analysis tools, or package-building systems.

They provide a representation of the workpiece state after a specified sequence of operations.

In-process objects are stored within the \texttt{in\_process} directory of the MIFX package.

Example location:

\begin{verbatim}
in_process/op1-5.json
\end{verbatim}

\subsection{In-Process Structure}

An In-Process object typically contains:

\begin{itemize}
\item an identifier
\item an optional descriptive name
\item optional coverage information
\item artifact references describing the workpiece state
\end{itemize}

Example:

\begin{verbatim}
{
  "id": "op1-5",
  "name": "After operations 1 to 5",

  "artifacts": [
    {
      "role": "in_process_mesh",
      "kind": "stl",
      "path": "meshes/op1-5.stl",
      "present": true
    }
  ]
}
\end{verbatim}

\subsection{Operation Coverage}

In-process objects typically correspond to a range of operations that produced the represented workpiece state.

The naming of the object may indicate the covered operation range.

The specification does not require a fixed naming convention.

\subsection{In-Process Artifacts}

In-process objects reference associated resources through artifacts.

Common artifacts include:

\begin{itemize}
\item workpiece geometry meshes
\item analysis metadata
\item optional extensions
\item optional addons
\end{itemize}

\subsection{Optional Nature}

In-process objects are optional within a MIFX package.

A valid MIFX file may omit the \texttt{in\_process} directory entirely.

Systems that do not support in-process data may safely ignore these objects.
